%------------------------------------------------------------------------------%
\begin{question}
Seja
\[
  f(x, y) = \ln\left(x^2 + y^2\right)
\]
use a aproximação linear de $f$ no ponto $(1, 2)$
para estimar $f(1,01; 1,98)$
\end{question}

%------------------------------------------------------------------------------%
\begin{resolution}{Avaliando a derivada parcial em $x$}
  \[
    \D{f}{x}
    \pause
    \;=\; \D{}{x} \left(\ln\left(x^2 + y^2\right)\right)
    \pause
    \;=\; \frac{1}{x^2 + y^2}\D{}{x} \left(x^2 + y^2\right)
    \pause
    \;=\; \frac{2x}{x^2 + y^2}
  \]

  \pause
  \bigskip
  \[
    \D{f}{x}(1, 2)
    \pause
    \;=\; \frac{2x}{x^2 + y^2} \evalat{(1, 2)}{}
    \pause
    \;=\; \frac{2\times 1}{1^2 + 2^2}
    \pause
    \;=\; \frac{2}{5}
  \]
\end{resolution}

%------------------------------------------------------------------------------%
\begin{resolution}{Avaliando a derivada parcial em $y$}
  \[
    \D{f}{y}
    \pause
    \;=\; \D{}{y} \left(\ln\left(x^2 + y^2\right)\right)
    \pause
    \;=\; \frac{1}{x^2 + y^2}\D{}{y} \left(x^2 + y^2\right)
    \pause
    \;=\; \frac{2y}{x^2 + y^2}
  \]

  \pause
  \bigskip
  \[
    \D{f}{y}(1, 2)
    \pause
    \;=\; \frac{2y}{x^2 + y^2} \evalat{(1, 2)}{}
    \pause
    \;=\; \frac{2\times 2}{1^2 + 2^2}
    \pause
    \;=\; \frac{4}{5}
  \]
\end{resolution}

%------------------------------------------------------------------------------%
\begin{resolution}{Avaliando a função}
  \[
    f(1, 2)
    \pause
    \;=\; \ln(1 + 4)
    \pause
    \;=\; \ln(5)
  \]
\end{resolution}

%------------------------------------------------------------------------------%
\begin{resolution}{Aproximação linear}
  \begin{align*}
    \onslide<+->{L(x, y) }
    \onslide<+->{& = f(a, b) + \D{f}{x}(a, b)(x-a) + \D{f}{y}(a, b)(y-b) \\[2mm]}
    \onslide<+->{& = \ln(5) + \frac{2}{5}(x - 1) + \frac{4}{5}(y - 2)}
  \end{align*}
\end{resolution}

%------------------------------------------------------------------------------%
\begin{resolution}{Avaliando no ponto $f(1,01; 1,98)$}
  \begin{align*}
    \onslide<+->{L(1,01; 1,98) }
    \onslide<+->{& = \ln(5) + \frac{2}{5}(0,01) + \frac{4}{5}(-0,02) \\}
    \onslide<+->{& = \ln(5) + \frac{2-8}{500} \\}
    \onslide<+->{& = \ln(5) - \frac{6}{500} \\}
    \onslide<+->{& = \ln(5) - \frac{3}{250}}
  \end{align*}
\end{resolution}

%------------------------------------------------------------------------------%
\begin{resolution}{Comparando com o valor exato}
  Com uma calculadora podemos avaliar
  \pause
  \[
    L(1,01; 1,98) \approx 1.59744
  \]
  \pause
  \[
    f(1,01; 1,98) \approx 1.59747
  \]
\end{resolution}
%------------------------------------------------------------------------------%
